\documentclass{article}
\usepackage[utf8]{inputenc}

\usepackage{amsmath, amssymb, amsthm, amsfonts}
\usepackage{subcaption, enumerate}
\usepackage{bbold}
\usepackage{polynom}
\usepackage{xfrac}
\usepackage{graphicx}
\usepackage{mathtools}
\usepackage{enumitem}

\usepackage{tikz}

\newtheorem{theorem}{Theorem}
\newtheorem{corot}{Corollary}[theorem]
\newtheorem{lemma}{Lemma}
\newtheorem{corol}{Corollary}[lemma]
\newtheorem{proposition}{Proposition}
\newtheorem{corop}{Corollary}[proposition]

\theoremstyle{definition}
\newtheorem{definition}{Definition}
\newtheorem{example}{Example}[definition]
\newtheorem*{notation}{Notation}

\newcommand{\dab}
{
	\begin{tikzpicture}[scale=0.04]
		% head
		\draw[thick] (0, 0) circle (3cm);

		% bent arm
		\draw[thick, rotate around={15 : (0, 0)}] (-3, 0) -- (-3.7, -2) -- (4.1, -1.1);

		% extended arm
		\draw[thick, rotate=20] (3, 0) -- (7, 2);
	\end{tikzpicture}
}
\newcommand{\st}{\text{ s.t. }}

\makeatletter
\renewcommand*\env@matrix[1][*\c@MaxMatrixCols c]{%
  \hskip -\arraycolsep
  \let\@ifnextchar\new@ifnextchar
  \array{#1}}
\makeatother

\usepackage{fullpage}

\title{Determinants Day 3 \\ with Dr Nevard}
\author{Laura Gallo}
\date{October 13, 2020}

\begin{document}
\maketitle

\section{Two theorems}
\begin{theorem}
	If $A$, $B$ are matrices, then $|AB|=|A||B|$.
	\begin{proof}
		Split the proof into two cases.
		\begin{description}
			\item[Case 1:] $|B|=0$ \\
				Then $B$ maps the domain into a smaller dimensional space, so $\text{dim Im}(AB)<\text{dim }AB$. \\
				This means the product is singular, so $|AB|=0=|A||B|$.
			\item[Case 2:] $|B|\neq 0$ \\
				If $|A|=0$, then the above argument can be used. \\
				Suppose $|A|\neq 0$. \\
				Let $f(A)=\frac{|AB|}{|B|}$. We claim that $f$ is multilinear, alternating, and normalised (ie, is the determinant by uniqueness).
				\begin{description}
					\item[Normalised:] $f(I)=\frac{|IB|}{|B|}=\frac{|B|}{|B|}=1$
					\item[Alternating] Swapping any two rows of $A$ is equivalent to swapping any two rows of $AB$ (left as an exercise to the reader). \\
						$f(A_\text{swap})=\frac{|A_\text{swap}B|}{|B|}=\frac{|(AB)_\text{swap}|}{|B|}=-\frac{|AB|}{|B|}=-f(A)$
					\item[Multilinear:] Let $v_1,\dots,av_k+bu_k,\dots,v_n$ be the rows of matrix $A$.
						\begin{equation*}
							f\left(
								\begin{matrix}
									v_1 \\
									v_2 \\
									\vdots \\
									av_k+bu_k \\
									\vdots \\
									v_n
								\end{matrix}
							\right) = \frac{1}{|B|}\left|
								\left(
									\begin{matrix}
										v_1 \\
										v_2 \\
										\vdots \\
										av_k+bu_k \\
										\vdots \\
										v_n
									\end{matrix}
								\right)B
							\right| = \frac{1}{|B|}\left|
								a\left(
									\begin{matrix}
										v_1 \\
										v_2 \\
										\vdots \\
										v_k \\
										\vdots \\
										v_n
									\end{matrix}
								\right)B+b\left(
									\begin{matrix}
										v_1 \\
										v_2 \\
										\vdots \\
										u_k \\
										\vdots \\
										v_n
									\end{matrix}
								\right)B
							\right|
						\end{equation*}
						\begin{equation*}
							= \frac{a}{|B|}\left|
								\left(
									\begin{matrix}
										v_1 \\
										v_2 \\
										\vdots \\
										v_k \\
										\vdots \\
										v_n
									\end{matrix}
								\right)B
							\right|+\frac{b}{|B|}\left|
								\left(
									\begin{matrix}
										v_1 \\
										v_2 \\
										\vdots \\
										u_k \\
										\vdots \\
										v_n
									\end{matrix}
								\right)B
							\right| = af\left(
								\begin{matrix}
									v_1 \\
									v_2 \\
									\vdots \\
									v_k \\
									\vdots \\
									v_n
								\end{matrix}
							\right) + bf\left(
								\begin{matrix}
									v_1 \\
									v_2 \\
									\vdots \\
									u_k \\
									\vdots \\
									v_n
								\end{matrix}
							\right)
						\end{equation*}
			\end{description}
		\end{description}
		Therefore by uniqueness, $|AB|=|A||B|$. \dab
	\end{proof}
\end{theorem}

\begin{theorem}
	Every matrix has a determinant.

	\begin{proof}
		By induction.
		\begin{description}
			\item[$n=1$:] Trivial.
			\item[$n=2$:] Already proven.
			\item[$n=k+1$:] Recall that for a diagonal matrix \begin{math}
				A=\left(
					\begin{matrix}
						a_{11} & 0 & \dots & 0 \\
						0 & a_{22} & \dots & 0 \\
						\vdots & \vdots & \ddots & \vdots \\
						0 & 0 & \dots & a_{nn}
					\end{matrix}
				\right)
			\end{math}, $|A|=\prod_{i=1}^na_{ii}$. \\
			We can permute the rows of diagonal matrices and sum them up with one element not from the main diagonal. For example:
			\begin{equation*}
				\left|
					\begin{matrix}
						0 & 0 & a & 0 \\
						b & 0 & 0 & 0 \\
						0 & c & 0 & 0 \\
						0 & 0 & 0 & d
					\end{matrix}
				\right| + \left|
					\begin{matrix}
						e & 0 & 0 & 0 \\
						b & 0 & 0 & 0 \\
						0 & c & 0 & 0 \\
						0 & 0 & 0 & d
					\end{matrix}
				\right| = \left|
					\begin{matrix}
						e & 0 & a & 0 \\
						b & 0 & 0 & 0 \\
						0 & c & 0 & 0 \\
						0 & 0 & 0 & d
					\end{matrix}
				\right|
			\end{equation*}
			From this, we claim that the determinant of a matrix is:
			$$|A|=\sum_{\pi\in S_n}(-1)^\pi \prod_{i=1}^n a_{i\pi_i}$$
			where $S_n$ is the group of permutations on a collection of $n$ items.

			\begin{description}
				\item[Normalisation:] Trivial.
				\item[Alternating:] Let $\sigma$ be the permutation of swapping two rows. \\
					$$\sum_{\pi\in S_n}(-1)^\pi\prod_{i=1}^na_{i(\pi\sigma)_i}=\sum_{\pi\in S_n}(-1)^{\pi\sigma}(-1)^\sigma\prod_{i=1}^na_{i(\pi\sigma)_i}=-\sum_{\pi'\in S_n}(-1)^{\pi'}\prod_{i=1}^na_{i\pi'_i}$$
				\item[Multilinear:] Let $a_{k\pi_k}=\alpha b_{\pi_k}+\beta c_{\pi_k}$
					$$\sum_{\pi\in S_n}(-1)^\pi\prod_{i=1}^na_{i\pi_i}=\sum_{\pi\in S_n}(-1)^\pi\left(\prod_{i=1}^{k-1}a_{i\pi_i}\right)a_{k\pi_k}\left(\prod_{i=k+1}^{n}a_{i\pi_i}\right)$$
					$$=\sum_{\pi\in S_n}(-1)^\pi\left(\prod_{i=1}^{k-1}a_{i\pi_i}\right)(\alpha b_{\pi_k}+\beta c_{\pi_k})\left(\prod_{i=k+1}^{n}a_{i\pi_i}\right)$$
					From here the proof is trivial. \dab
			\end{description}
		\end{description}
	\end{proof}
\end{theorem}

\end{document}
